\section{正道少年的荒蛮历险：从 App Store 到 NLP}

前面讲最后访谈的特殊性，这一章回到 Peak 的起点。Peak 的早期经历很像一条微型技术史。他在高中时期做第三方 iOS 浏览器，赶上 App Store 早期红利，用最朴素的 buy-copy 模式获得现金流。这个经历让他很早理解：平台变化会让个人开发者、创业公司和大厂短暂站到同一起跑线。

\begin{figure}[H]
\centering
\includegraphics[width=0.96\textwidth]{manus-journey.png}
\caption{Manus 奇幻漂流路线：从 App Store、NLP、Open IE 到 Manus 和 Meta 收购。}
\end{figure}

\begin{knowledgebox}{读图：Manus 不是从天而降}
图里每个节点都来自一个具体需求。浏览器需要弱网下更快，于是研究下一次点击预测；点击预测引向 NLP；NLP 引向 Word2Vec、知识图谱和 Open IE；这些经验最终进入 Agent 产品化。
\end{knowledgebox}

\subsection{App Store 蛮荒期}

本节先看第一个平台窗口。Peak 说，App Store 给了他一个第三方指标：自己瞎捣鼓的软件能产生经济价值。这个指标对一个高中生非常重要，因为它让“不务正业”变成可验证的产品能力。移动互联网早期的新媒介、新分发和新商业模式，给个人开发者留下了窗口。

\begin{figure}[H]
\centering
\includegraphics[width=0.94\textwidth]{appstore-to-ai-platform.png}
\caption{App Store 蛮荒期与 AI 时代：平台红利的差异。}
\end{figure}

\begin{knowledgebox}{读图：AI 技术突破不等于新平台红利}
左侧 App Store 蛮荒期带来硬件媒介变化和分发红利；右侧 AI 时代虽然技术突破巨大，但没有同等清晰的新平台，所以巨头、创业公司和个人开发者反应都很快，窗口更窄。
\end{knowledgebox}

\subsection{从浏览器预加载到 NLP}

浏览器时代，他想预测用户下一次点击，提前加载下一页，从而改善弱网体验。这个具体需求把他带进 NLP。后来 Word2Vec 让他看到自然语言可以被转成稠密向量，知识图谱则让他尝试用结构化方式回答用户问题。

\begin{figure}[H]
\centering
\includegraphics[width=0.96\textwidth]{nlp-lineage.png}
\caption{从浏览器预加载到 NLP：具体需求牵引学习路径。}
\end{figure}

\begin{knowledgebox}{读图：学习路径由问题牵引}
图中从浏览器到 NLP，不是先选赛道再找问题，而是从“弱网时怎么更快”这个具体需求出发。好的技术学习往往不是抽象追热点，而是被真实产品问题牵引。
\end{knowledgebox}

\subsection{为什么浏览器不适合创业者长期做}

Peak 也反思过，第三方浏览器从古至今都不太适合创业者长期做。浏览器更像巨头有分发渠道后做的入口型产品，因为它天然要求系统级分发、默认入口、生态协作和安全信任。这个判断对今天的 AI 浏览器也有启发：AI 浏览器听起来像创业机会，但如果没有分发和系统入口，创业公司很容易被巨头挤压。

Manus 与浏览器的关系更微妙。它不是简单做一个浏览器，而是试图让 Agent 在浏览器和网页世界里行动。也就是说，它不一定要成为浏览器本身，而是要成为跨网页、工具和任务的行动层。这解释了为什么它既继承了 Peak 早期浏览器/NLP 的问题意识，又没有回到“做浏览器”这条路。

\begin{longtable}{p{0.24\textwidth}p{0.34\textwidth}p{0.33\textwidth}}
\toprule
对象 & 核心任务 & 创业难点 \\
\midrule
传统浏览器 & 呈现网页、管理标签、承载搜索和网站入口 & 默认入口和分发权通常被巨头掌握。 \\
AI 浏览器 & 在浏览器里加入总结、问答、自动填写等能力 & 如果只是浏览器增强，仍受入口和生态约束。 \\
Agent 行动层 & 跨网页、工具、文件和账户完成任务 & 需要权限、支付、反爬协作和可靠执行。 \\
\bottomrule
\end{longtable}

\begin{importantbox}{从浏览器入口到 Agent 行动层}
浏览器负责呈现网页，Agent 负责理解目标并行动。AI 浏览器如果只是“更聪明的浏览器”，可能仍受入口限制；Agent 如果能跨工具执行任务，就可能成为浏览器之上的行动层。
\end{importantbox}

\subsection{本章小结}

Manus 的底层气质来自早期个人开发者经历：先解决具体问题，再沿着问题学习技术。这个路径解释了 Peak 为什么更重视产品中的真实反馈，而不是只讲宏大 AI 叙事。

